%==============================================================================
% COMANDOS PERSONALIZADOS
% Livro: Odontologia & Inteligência Artificial
%==============================================================================

% --- Compatibilidade biblatex: \citeonline → \textcite ---
\providecommand{\citeonline}{\textcite}

%-----------------------------------------------------------------------------
% 1. NÍVEIS DE DIFICULDADE (bolinhas coloridas)
%-----------------------------------------------------------------------------
% Uso: \nivelzero, \nivelum, \niveldois, \niveltres, \nivelquatro
% Cada comando insere bolinhas preenchidas/vazias indicando o nível.

\newcommand{\nivelzero}{%
  \textcolor{corNivelFacil}{%
    \raisebox{-0.2ex}{\tikz{\fill (0,0) circle (2.5pt);}}%
    \hspace{2pt}\textit{Nível 0 — Iniciante Absoluto}%
  }%
}

\newcommand{\nivelbasico}{%
  \textcolor{corNivelFacil}{%
    \raisebox{-0.2ex}{\tikz{\fill (0,0) circle (2.5pt);}}%
  }%
  \textcolor{corNivelFacil}{\raisebox{-0.2ex}{\tikz{\fill (4pt,0) circle (2.5pt);}}}%
  \hspace{2pt}\textit{Nível Básico}%
}

\newcommand{\nivelum}{%
  \textcolor{corNivelFacil}{\raisebox{-0.2ex}{\tikz{\fill (0,0) circle (2.5pt);}}}%
  \textcolor{lightgray}{\raisebox{-0.2ex}{\tikz{\fill[lightgray] (4pt,0) circle (2.5pt);}}}%
  \textcolor{lightgray}{\raisebox{-0.2ex}{\tikz{\fill[lightgray] (8pt,0) circle (2.5pt);}}}%
  \textcolor{lightgray}{\raisebox{-0.2ex}{\tikz{\fill[lightgray] (12pt,0) circle (2.5pt);}}}%
  \textcolor{lightgray}{\raisebox{-0.2ex}{\tikz{\fill[lightgray] (16pt,0) circle (2.5pt);}}}%
  \hspace{2pt}\textit{Nível Iniciante}%
}

\newcommand{\niveldois}{%
  \textcolor{corNivelMedio}{\raisebox{-0.2ex}{\tikz{\fill (0,0) circle (2.5pt);}}}%
  \textcolor{corNivelMedio}{\raisebox{-0.2ex}{\tikz{\fill (4pt,0) circle (2.5pt);}}}%
  \textcolor{lightgray}{\raisebox{-0.2ex}{\tikz{\fill[lightgray] (8pt,0) circle (2.5pt);}}}%
  \textcolor{lightgray}{\raisebox{-0.2ex}{\tikz{\fill[lightgray] (12pt,0) circle (2.5pt);}}}%
  \textcolor{lightgray}{\raisebox{-0.2ex}{\tikz{\fill[lightgray] (16pt,0) circle (2.5pt);}}}%
  \hspace{2pt}\textit{Nível Intermediário}%
}

\newcommand{\niveltres}{%
  \textcolor{corNivelDificil}{\raisebox{-0.2ex}{\tikz{\fill (0,0) circle (2.5pt);}}}%
  \textcolor{corNivelDificil}{\raisebox{-0.2ex}{\tikz{\fill (4pt,0) circle (2.5pt);}}}%
  \textcolor{corNivelDificil}{\raisebox{-0.2ex}{\tikz{\fill (8pt,0) circle (2.5pt);}}}%
  \textcolor{lightgray}{\raisebox{-0.2ex}{\tikz{\fill[lightgray] (12pt,0) circle (2.5pt);}}}%
  \textcolor{lightgray}{\raisebox{-0.2ex}{\tikz{\fill[lightgray] (16pt,0) circle (2.5pt);}}}%
  \hspace{2pt}\textit{Nível Avançado}%
}

\newcommand{\nivelquatro}{%
  \textcolor{corNivelAvancado}{\raisebox{-0.2ex}{\tikz{\fill (0,0) circle (2.5pt);}}}%
  \textcolor{corNivelAvancado}{\raisebox{-0.2ex}{\tikz{\fill (4pt,0) circle (2.5pt);}}}%
  \textcolor{corNivelAvancado}{\raisebox{-0.2ex}{\tikz{\fill (8pt,0) circle (2.5pt);}}}%
  \textcolor{corNivelAvancado}{\raisebox{-0.2ex}{\tikz{\fill (12pt,0) circle (2.5pt);}}}%
  \textcolor{lightgray}{\raisebox{-0.2ex}{\tikz{\fill[lightgray] (16pt,0) circle (2.5pt);}}}%
  \hspace{2pt}\textit{Nível Especialista}%
}

\newcommand{\nivelcinco}{%
  \textcolor{corNivelPhD}{\raisebox{-0.2ex}{\tikz{\fill (0,0) circle (2.5pt);}}}%
  \textcolor{corNivelPhD}{\raisebox{-0.2ex}{\tikz{\fill (4pt,0) circle (2.5pt);}}}%
  \textcolor{corNivelPhD}{\raisebox{-0.2ex}{\tikz{\fill (8pt,0) circle (2.5pt);}}}%
  \textcolor{corNivelPhD}{\raisebox{-0.2ex}{\tikz{\fill (12pt,0) circle (2.5pt);}}}%
  \textcolor{corNivelPhD}{\raisebox{-0.2ex}{\tikz{\fill (16pt,0) circle (2.5pt);}}}%
  \hspace{2pt}\textit{Nível PhD}%
}

%-----------------------------------------------------------------------------
% 2. NOTA EXPLICATIVA (nota de rodapé especial com marcador próprio)
%-----------------------------------------------------------------------------
% Uso: \notaexplicativa{Texto da nota explicativa aqui.}

\newcommand{\notaexplicativa}[1]{%
  \footnote{\textbf{Nota explicativa:} #1}%
}

%-----------------------------------------------------------------------------
% 3. INSERÇÃO DE CÓDIGO (listings)
%-----------------------------------------------------------------------------
% Uso: \codigo{caminho/do/arquivo.py}
% Insere o conteúdo de um arquivo Python com a formatação listings.

\newcommand{\codigo}[1]{%
  \lstinputlisting[language=Python]{#1}%
}

%-----------------------------------------------------------------------------
% 4. RESULTADO DE TESTE (TDD)
%-----------------------------------------------------------------------------
% Uso: \teste{descrição do teste}
% Cria uma caixa destacada com o resultado do teste.

\newcommand{\teste}[1]{%
  \begin{mdframed}[
    linecolor=corTesteSucesso,
    linewidth=2pt,
    backgroundcolor=corTesteSucesso!5,
    roundcorner=4pt,
    skipabove=10pt,
    skipbelow=10pt
  ]
    \textbf{\textcolor{corTesteSucesso}{\checkmark\ Teste TDD:}} #1
  \end{mdframed}%
}

\newcommand{\testefalha}[1]{%
  \begin{mdframed}[
    linecolor=corTesteFalha,
    linewidth=2pt,
    backgroundcolor=corTesteFalha!5,
    roundcorner=4pt,
    skipabove=10pt,
    skipbelow=10pt
  ]
    \textbf{\textcolor{corTesteFalha}{$\times$\ Teste TDD FALHOU:}} #1
  \end{mdframed}%
}

%-----------------------------------------------------------------------------
% 5. CITAÇÃO ATIVA COM DOI (ABNT NBR 10520)
%-----------------------------------------------------------------------------
% Uso: \citacaoativa{doi-xxxx-xxxxx}{10.xxxx/xxxxx}{Texto citado}
% Cita um trecho longo em bloco (recuado, sem aspas) com referência ABNT.
% O primeiro argumento é a chave no .bib (para gerar entrada na bibliografia).
% O segundo é o DOI (exibido como link).
% O terceiro é o texto citado (sem aspas, conforme NBR 10520).

\newcommand{\citacaoativa}[3]{%
  \cite{#1}%
  \begin{citacao}
    #3%
    \par\hfill--- (\citeonline{#1}, DOI: \url{https://doi.org/#2})
  \end{citacao}%
}

% Comando alternativo para citação sem chave .bib (não gera entrada na bibliografia)
\newcommand{\citacaoativasimples}[2]{%
  \begin{citacao}
    #2%
    \par\hfill--- (DOI: \url{https://doi.org/#1})
  \end{citacao}%
}

%-----------------------------------------------------------------------------
% 6. BLOCO SDD (Specification-Driven Development)
%-----------------------------------------------------------------------------
% Uso: \spec{Texto da especificação aqui.}

\newcommand{\spec}[1]{%
  \begin{mdframed}[
    linecolor=corPrimaria,
    linewidth=2pt,
    backgroundcolor=corPrimaria!5,
    roundcorner=4pt,
    leftmargin=10pt,
    rightmargin=10pt,
    innertopmargin=10pt,
    innerbottommargin=10pt
  ]
    \textbf{\textcolor{corPrimaria}{SDD — Especificação:}}\par
    #1
  \end{mdframed}%
}

%-----------------------------------------------------------------------------
% 7. FIGURA DE ARTIGO REAL
%-----------------------------------------------------------------------------
% Uso: \artigofig{caminho/imagem}{doi_referencia}{Legenda descritiva}
% Cria uma figura com referência a artigo real.

\newcommand{\artigofig}[3]{%
  \begin{figure}[H]
    \centering
    \includegraphics[width=0.85\textwidth]{#1}
    \caption[#3]{#3 \\(Fonte: \citeonline{#2})}
    \label{fig:#1}
  \end{figure}%
}

%-----------------------------------------------------------------------------
% 8. CAIXA DE DESTAQUE (informação importante)
%-----------------------------------------------------------------------------
% Uso: \destaque{Texto em destaque}

\newcommand{\destaque}[1]{%
  \begin{mdframed}[
    linecolor=corDestaque,
    linewidth=2pt,
    backgroundcolor=corDestaque!5,
    roundcorner=4pt,
    leftmargin=10pt,
    rightmargin=10pt,
    innertopmargin=10pt,
    innerbottommargin=10pt
  ]
    \textbf{\textcolor{corDestaque}{Destaque:}}\par #1
  \end{mdframed}%
}

%-----------------------------------------------------------------------------
% 9. AMBIENTE DE EXEMPLO PRÁTICO
%-----------------------------------------------------------------------------
% Uso: \begin{pratica}{Título da prática} ... \end{pratica}

\newenvironment{pratica}[1][]{%
  \begin{mdframed}[
    linecolor=corSecundaria,
    linewidth=2pt,
    backgroundcolor=corSecundaria!5,
    roundcorner=4pt,
    leftmargin=10pt,
    rightmargin=10pt,
    innertopmargin=10pt,
    innerbottommargin=10pt
  ]
    \textbf{\textcolor{corSecundaria}{Prática: #1}}\par
}{%
  \end{mdframed}%
}

%-----------------------------------------------------------------------------
% 10. COMANDO PARA GLOSSÁRIO
%-----------------------------------------------------------------------------
% Termos técnicos com entrada no glossário

\newcommand{\termo}[2]{%
  \textbf{#1}\gls{#2}%
}

%-----------------------------------------------------------------------------
% 11. COMANDO PARA ABREVIAÇÕES
%-----------------------------------------------------------------------------
% Uso: \sigla{IA}{Inteligência Artificial}

\newcommand{\sigla}[2]{%
  \acronym{#1}{#2}%
}

%-----------------------------------------------------------------------------
% 12. ENTRADA NO ÍNDICE REMISSIVO
%-----------------------------------------------------------------------------
% Uso: \indice{termo}
% Atalho para \index

\newcommand{\indice}[1]{\index{#1}}

%-----------------------------------------------------------------------------
% 13. VERSÃO DO LIVRO
%-----------------------------------------------------------------------------
\newcommand{\versaolivro}{v1.0.0}
\newcommand{\dataversao}{\today}